\documentclass[../../../include/open-logic-section]{subfiles}

\begin{document}

\olfileid{sfr}{infinite}{card-sb}

\olsection{ضميمه: د شروډر--برنشتاين ثبوت}

له ساده سټ نظريې څخه د وتلو مخکښې يو وروستے ساده (خو ډېر فني!)
ثبوت څېړو۔ دا د شروډر--برنشتاين ثبوت دے
(\olref[sfr][siz][sb]{thm:schroder-bernstein}): که
$\cardle{A}{B}$ او $\cardle{B}{A}$ وي، نو $\cardeq{A}{B}$؛ يعنې که
!!{injection}s $f \colon A \to B$ او $g \colon B \to A$ راکړل
شي، نو يوه !!a{bijection} $h \colon A \to B$ شته۔

په دې څپرکي کښې مو د ډېډېکېنډ د \emph{بندښتونو} مفهوم وکاراوه۔ په
رښتيا، ډېډېکېنډ په همدې مفهوم د شروډر--برنشتاين يو ښکلی ثبوت ورکړ،
چې دلته ئې وړاندې کوو۔ که لږ بېله خو په اصل کښې ورته شننه غواړئ،
دا ثبوت په نژدې ډول د \citet[pp.~157--8]{Potter2004} پيروي کوي۔
په انټرنېټ کښې لږه پلټنه به هم درته وښيي چې دا---لکه د
$\sqrt{2}$ غير ناطق‌والے---هغه قضيه ده چې \emph{ډېر} په زړه پورې
او بېل ثبوتونه لري۔

د \olref[sfr][infinite][dedekind]{Closure} په شان نښه وکاروئ او د
هر سټ $B$ او تابعې~$f$ دپاره وټاکئ:
\[
\Closureofunder{f}{B} = \bigcap \Setabs{X}{B \subseteq X
\text{ او $X$ د $f$ له مخې بند دے}}
\]
په دې ډول تعريف شوے $\Closureofunder{f}{B}$ تر ټولو کوچنے داسې د
$f$ له مخې بند سټ دے چې~$B$ رانغاړي، يعنې:

\begin{lem}\ollabel{Closureprops}
	د هرې تابع $f$ او هر $B$ دپاره:
	\begin{enumerate}
		\item\ollabel{Closurehaselem} $B \subseteq \Closureofunder{f}{B}$؛ او
		\item\ollabel{Closureclosed} $\Closureofunder{f}{B}$ د $f$ له مخې بند دے؛ او
		\item\ollabel{Closuresmallest} که $X$ د $f$ له مخې بند وي او $B
		\subseteq X$ وي، نو $\Closureofunder{f}{B} \subseteq X$۔
	\end{enumerate}
\end{lem}

\begin{proof}
کټ مټ لکه په \olref[sfr][infinite][dedekind]{closureproperties} کښې۔
\end{proof}

شروډر--برنشتاين ته د رسېدو دپاره يو وروستے حقيقت پکار دے:

\begin{prop}\ollabel{sbhelper}
که $A \subseteq B \subseteq C$ او $A \approx C$ وي، نو $\cardeq{A}{B}$۔
\end{prop}
\textbf{د ژباړې د نسخې سمون (OLINF-002):} د سرچينې په پايله کښې
د کارډينال برابرۍ يوه نښه په تېروتنه د بلې کارډينال برابرۍ د لومړي
ارګومېنټ په توګه ځالې شوې ده۔ لاندې ثبوت له کوچني سټ څخه منځني سټ
ته دوه اړخيزه يو پر يو تابع جوړوي، نو سم بيان هماغه د دغو دوو سټونو
همشمېر والی دے چې پورته ليکل شوے۔

\begin{proof}
يوه !!a{bijection} $f \colon C \to A$ راکړل شي، وټاکئ
$F = \Closureofunder{f}{C \setminus B}$ او له ډومېن $C$ سره تابع
$g$ داسې تعريف کړئ:
\[
	g(x) =
	\begin{cases}
			f(x) &\text{که $x \in F$ وي}\\
			x & \text{که نۀ}
	\end{cases}
\]
ښيو چې $g$ له $C \to B$ يوه !!a{bijection} ده؛ له دې به راووځي
چې $\comp{f^{-1}}{g} \colon A \to B$ يوه !!a{bijection} ده، او
ثبوت بشپړېږي۔

لومړے ادعا کوو چې که $x \in F$ خو $y\notin F$ وي، نو
$g(x) \neq g(y)$۔ د نقيض دپاره فرض کړئ چې برابر دي؛ نو
$y = g(y) = g(x) = f(x)$۔ ځکه $x \in F$ او د
\olref{Closureprops} له مخې $F$ د $f$ له مخې بند دے،
$y = f(x) \in F$ لرو، چې ټکر دے۔

اوس فرض کړئ $g(x) = g(y)$۔ د پورته خبرې له مخې، $x \in F$ هله او
يوازې هله چې $y \in F$ وي۔ که $x, y \in F$ وي، نو
$f(x) = g (x) = g(y) = f(y)$؛ او ځکه $f$ يوه !!a{bijection} ده،
$x = y$۔ که $x, y \notin F$ وي، نو $x = g(x) = g(y) = y$۔ نو
$g$ يوه !!a{injection} ده۔

دا پاتې ده چې وښيو $\ran{g} = B$۔ يو $x \in B \subseteq C$
وټاکئ۔ که $x \notin F$ وي، نو $g(x) = x$۔ که $x \in F$ وي، نو د
کوم $y \in F$ دپاره $x = f(y)$؛ ځکه که داسې نۀ وي، نو
$F \setminus \{x\}$ به د $f$ له مخې بند وي او $C\setminus B$ به
وغځوي، چې د \olref{Closureprops} له مخې ناشونې ده؛ اوس
$g(y) = f(y) = x$۔
\end{proof}

په پاے کښې د اصلي پايلې ثبوت دا دے۔ ياد کړئ چې د تابعې $h$ او سټ
$D$ دپاره $\funimage{h}{D} = \Setabs{h(x)}{x \in D}$ تعريفوو۔

\begin{proof}[د شروډر--برنشتاين ثبوت] فرض کړئ $f \colon A \to B$
او $g \colon B \to A$ !!{injection}s دي۔ ځکه
$\funimage{f}{A} \subseteq B$، نو
$\funimage{g}{\funimage{f}{A}} \subseteq g[B] \subseteq A$۔
برسېره پر دې، $\comp{f}{g} \colon A \to
\funimage{g}{\funimage{f}{A}}$ يوه !!{injection} ده، ځکه $g$ او
$f$ دواړه همداسې دي؛ او په رښتيا $\comp{f}{g}$ يوه
!!a{bijection} ده، يوازې په دې سبب چې مقابل ډومېن مو همداسې تعريف
کړے۔ نو $\cardeq{\funimage{g}{\funimage{f}{A}}}{A}$، او په دې توګه
د \olref{sbhelper} له مخې يوه !!a{bijection} $h \colon A \to
\funimage{g}{B}$ شته۔ برسېره پر دې، $g^{-1}$ يوه !!a{bijection}
$\funimage{g}{B} \to B$ ده۔ نو $\comp{h}{g^{-1}} \colon A \to B$
يوه !!a{bijection} ده۔
\end{proof}

\end{document}
